\section{Introduction}
\label{sec:intro}

Pretraining on large-scale vision and vision-language data has substantially improved object detection~\cite{detr,dino,groundingdino}. Nevertheless, detectors trained on natural images remain brittle in aerial, industrial, underwater, and artistic domains, which differ in viewpoint, scale, texture, illumination, and scene composition. Because dense annotation in these domains is expensive, cross-domain few-shot object detection (CD-FSOD) seeks to adapt a source-pretrained detector using only a few labeled examples per category~\cite{cdvito,domainrag}.

Recent work employs feature caching and sparse adaptation~\cite{ifc,sivito}, augmentation and retrieval-guided generation~\cite{ets,domainrag}, foundation-model experts~\cite{vfmmoe,yu2026}, or multimodal prototypes and visual prompts~\cite{lmp,petdino,astigmatism}. These methods share a common premise: scarce target examples contain useful domain evidence. This premise, however, is incomplete. A support image contains category identity, relevant foreground appearance, and incidental scene context. A DIOR airplane, for example, provides a useful near-nadir silhouette and small-object texture, but may also appear alongside runways or buildings that do not define the category. A mixed prototype can improve domain alignment while allowing such incidental context to shift the pretrained category boundary.

We investigate this failure mode through the controlled intervention shown in Fig.~\ref{fig:context-leakage}. We preserve the support foreground, bounding box, scale, and query set, changing only the pixels outside a protected support region. Across Clipart1k and NEU-DET, this intervention alters category distributions and predicted boxes beyond the variation induced by training seeds. We refer to this phenomenon as \emph{support-induced context leakage}. Its severity is domain dependent: neutralization is the most disruptive intervention for several NEU-DET diagnostics, whereas mismatched context causes the clearest aggregate degradation on Clipart1k. Context sensitivity therefore reflects the interaction among support evidence, domain statistics, and detector adaptation, rather than a fixed ranking of corruption severity.

This observation shifts the adaptation question from \emph{how much target knowledge can be injected?} to \emph{when should target knowledge be allowed to alter a decision?} We address this question with the knowledge-enhanced Hard-Soft framework illustrated in Fig.~\ref{fig:overview}. Foreground-Frozen Context Probing (FFCP) estimates where and to what extent context enters each category representation. Context-Constrained Hard-Soft Decomposition (CHSD) constructs a Hard identity anchor and a reliability-weighted Soft target-domain prototype after removing context-sensitive directions from the original support foreground. Asymmetric Task-Authority Routing (ATAR) uses agreement between these two estimates to determine prototype reliability, then applies a bounded cosine correction only to uncertain Grounding DINO queries. The framework retains a single detector and preserves its original localization path.

We evaluate the resulting method across six target domains and three shot regimes. Beyond aggregate detection accuracy, we measure fixed-query context sensitivity and inspect whether improvements arise from reduced background responses or from indiscriminate score changes. The method does not discard target-domain evidence; it preserves useful foreground variation while restricting unsupported category corrections.

Our contributions are summarized as follows:
\begin{itemize}
    \item We identify support-induced context leakage and introduce a foreground-preserving, fixed-query protocol that separates context effects from ordinary training-seed variation, verifying the phenomenon in both artistic and industrial target domains.
    \item We propose FFCP and CHSD to estimate context-sensitive directions and sequentially decompose support evidence into bounded Hard identity knowledge and context-cleaned Soft appearance knowledge.
    \item We introduce ATAR, a reliability-constrained Grounding DINO interface that uses Hard--Soft agreement and query uncertainty to bound class-score correction while preserving the baseline localization path. We evaluate its accuracy and context robustness across six CD-FSOD benchmarks.
\end{itemize}
